% ===================== Chapter 7 =====================
\chapter{Arrays and Vectors}
\label{ch:arrays-vectors}

A single variable holds one value. But programs constantly deal with
\emph{collections} a row of scores, a list of names, a grid of cells. Salam
gives you two tools for holding many values under one name: the \textbf{array},
whose size is fixed when you write it, and the \textbf{vector}, which grows and
shrinks as your program runs. This chapter covers both.

\section{Fixed-size arrays}

An \textbf{array} is a numbered row of values, all of the same type. You declare
one by writing the element type followed by the length in square brackets, and
you can give the initial contents in a bracketed list:

\begin{salam}
func main:
    scores i32[5] = [90, 85, 70, 95, 60]
    println scores[0]      // the first element
    println scores[4]      // the last element
end
\end{salam}

\begin{progout}
90
60
\end{progout}

The type \code{i32[5]} reads ``an array of five \code{i32}s.'' The values are
stored in \emph{slots} numbered from \code{0}, so a length-5 array has slots
\code{0}, \code{1}, \code{2}, \code{3}, \code{4}. You reach a slot with square
brackets: \code{scores[0]} is the first, \code{scores[4]} the last.

\begin{warn}
Indices start at \emph{zero}, not one. The first element is \code{scores[0]} and
the \emph{last} of a length-$n$ array is \code{scores[n-1]} here
\code{scores[4]}, not \code{scores[5]}. Off-by-one slips around this rule are
among the most common bugs in all of programming; the zero-based habit becomes
second nature with practice.
\end{warn}

\subsection{Changing elements}

To modify an element, assign to its slot. As with any modification, the array
must be \code{mut}:

\begin{salam}
func main:
    mut a i32[3] = [1, 2, 3]
    a[1] = 20
    println a[0], a[1], a[2]
end
\end{salam}

\begin{progout}
1 20 3
\end{progout}

\subsection{Walking an array with a loop}

The \code{for} loop and arrays are made for each other: the loop variable steps
through the valid indices, and the body does something with each element. Here we
sum an array and find its largest value in a single pass:

\begin{salam}
func main:
    a i32[7] = [5, 9, 1, 7, 3, 8, 2]
    mut sum i32 = 0
    mut hi i32 = a[0]
    for mut i = 0; i < 7; i = i + 1:
        sum += a[i]
        if a[i] > hi: hi = a[i] end
    end
    println "sum =", sum, "max =", hi
end
\end{salam}

\begin{progout}
sum = 35 max = 9
\end{progout}

Notice the loop condition \code{i < 7}: it stops \emph{before} index 7, visiting
exactly slots 0 through 6. This half-open pattern start at 0, continue while
\code{i < length} is the canonical way to walk an array, and it lines up
perfectly with zero-based indexing.

\section{Multi-dimensional arrays}

An array's elements can themselves be arrays, giving you a grid. \code{i32[2][3]}
is ``2 rows of 3 columns.'' You index with two pairs of brackets, row first:

\begin{salam}
func main:
    mut grid i32[2][3] = [[1, 2, 3], [4, 5, 6]]
    grid[1][2] = 60
    println grid[0][0], grid[1][2]
end
\end{salam}

\begin{progout}
1 60
\end{progout}

Nested loops walk a grid naturally: an outer loop over rows, an inner loop over
columns. This is the foundation for anything tabular game boards, matrices,
spreadsheets.

\section{Slices: a window onto an array}

Often you want to work with just \emph{part} of an array the middle four
elements, everything after the third, the first half without copying it into a
new array. A \textbf{slice} is a lightweight \emph{view} onto a stretch of an
array. It does not own any data and it does not allocate memory on the heap;
under the hood it is nothing more than a pointer to the first element plus a
length. Languages such as Go and Rust have the same idea.

You create a slice with a \textbf{range} inside the brackets, \code{a[lo:hi]}:

\begin{salam}
func main:
    nums i32[6] = [10, 20, 30, 40, 50, 60]
    middle auto = nums[1:4]
    println middle[0], middle[1], middle[2]
    println len(middle)
end
\end{salam}

\begin{progout}
20 30 40
3
\end{progout}

The range is \textbf{half-open}: \code{nums[1:4]} starts \emph{at} index~1 and
stops \emph{before} index~4, so it covers indices 1, 2 and 3 a length of
$4-1=3$. This is the same half-open convention as the loop condition
\code{i < length}, and it has a tidy consequence: the length of \code{a[lo:hi]}
is simply \code{hi - lo}, and \code{a[lo:hi]} followed by \code{a[hi:end]} covers
the whole range with no gap and no overlap.

\subsection{Leaving out a bound}

If you omit the low bound it defaults to \code{0}; omit the high bound and it runs
to the end of the array. Omit both, \code{a[:]}, and you get a view of the whole
array.

\begin{salam}
func main:
    nums i32[6] = [10, 20, 30, 40, 50, 60]
    println len(nums[:3])     // first three: 10, 20, 30
    println len(nums[3:])     // from index 3 on: 40, 50, 60
    println len(nums[:])      // the whole array
end
\end{salam}

\begin{progout}
3
3
6
\end{progout}

\subsection{The slice type \texttt{T[:]}}

A slice has its own type, written \code{T[:]} a slice of \code{T} as opposed to
\code{T[N]}, a fixed array of \code{N} items. You will most often see it as a
parameter type, because a function that takes \code{i32[:]} can be handed
\emph{any} run of \code{i32}s: a whole array, or just a piece of one.

\begin{salam}
func total(view i32[:]) i32:
    mut sum i32 = 0
    for mut i = 0; i < len(view); i = i + 1:
        sum = sum + view[i]
    end
    ret sum
end

func main:
    nums i32[6] = [10, 20, 30, 40, 50, 60]
    println total(nums[:])      // all six
    println total(nums[2:5])    // just 30 + 40 + 50
end
\end{salam}

\begin{progout}
210
120
\end{progout}

One \code{total} function now works on the whole array and on any slice of it. You
do not need a separate copy, and you do not pass the length as an extra argument
\code{len} reads it straight from the slice.

\subsection{A slice is a window, not a copy}

Because a slice borrows the array's storage rather than copying it, writing
through a slice changes the original array. This is the key difference from
copying out a sub-array, and it is exactly what makes slices cheap.

\begin{salam}
func main:
    mut nums i32[6] = [10, 20, 30, 40, 50, 60]
    mut tail auto = nums[4:]
    tail[0] = 500
    tail[1] = 600
    println nums[4], nums[5]
end
\end{salam}

\begin{progout}
500 600
\end{progout}

Assigning to \code{tail[0]} reaches through the window and overwrites
\code{nums[4]}. The slice stored no data of its own; it only ever pointed at the
array.

\begin{warn}
A slice borrows storage it does not own. Keep the underlying array alive for as
long as you use the slice, and remember that a slice has a \emph{fixed} length:
like a fixed array it does not grow. When you need something that grows, reach for
a vector (next).
\end{warn}

\section{The limit of fixed arrays}

A fixed array's length is part of its type and is decided when you write the
program. That is perfect when you know the count up front (a week has 7 days, a
board has 64 squares). But often you do \emph{not} know how many items there will
be you are reading lines from a file, or collecting whatever the user types.
For that you need a collection that can grow. Enter the vector.

\section{Vectors: arrays that grow}

A \textbf{vector}, written \code{Vector<T>}, is a sequence of \code{T} values
whose length can change at run time. You start with an empty one and add to it as
you go. Create an empty vector with \code{Vector \{\}}:

\begin{salam}
func main:
    mut v Vector<i32> = Vector {}
    v.push(10)
    v.push(20)
    v.push(30)
    println "length:", v.len()
end
\end{salam}

\begin{progout}
length: 3
\end{progout}

The angle brackets in \code{Vector<i32>} say what the vector holds here,
\code{i32}s. (This is your first glimpse of \emph{generics}, the subject of
Chapter~12: \code{Vector} is one type that works for any element type.) The
\code{.push(x)} method appends a value to the end, and \code{.len()} reports how
many values are currently stored.

\subsection{Reading and writing elements}

A vector's \code{.get(i)} method returns a \emph{pointer} to the element at index
\code{i}; to read the value itself, follow the pointer with \code{[0]}. To
overwrite an element, use \code{.set(i, value)}:

\begin{salam}
func main:
    mut v Vector<i32> = Vector {}
    v.push(10)
    v.push(20)
    v.push(30)
    println "first:", v.get(0)[0]
    v.set(0, 100)
    println "first now:", v.get(0)[0]
end
\end{salam}

\begin{progout}
first: 10
first now: 100
\end{progout}

\begin{note}
Why \code{v.get(i)[0]} rather than just \code{v[i]}? \code{.get} returns a
\emph{pointer} into the vector's storage, and \code{[0]} reads through that
pointer to the value. The pointer form is what makes vectors efficient no
copying and you will understand it fully after the chapter on memory
(Chapter~18). For now, treat \code{v.get(i)[0]} as the idiom for ``the value at
index \code{i}.''
\end{note}

\subsection{Walking a vector}

Because \code{.len()} can change, you typically walk a vector with a \code{while}
loop bounded by its current length:

\begin{salam}
func main:
    mut v Vector<i32> = Vector {}
    v.push(5)
    v.push(8)
    v.push(2)
    mut total i32 = 0
    mut i i32 = 0
    while i < v.len():
        total += v.get(i)[0]
        i += 1
    end
    println "total:", total
end
\end{salam}

\begin{progout}
total: 15
\end{progout}

\subsection{Removing from the end}

\code{.pop()} removes the last element and returns it turning a vector into a
\emph{stack}, where the last thing in is the first thing out:

\begin{salam}
func main:
    mut v Vector<i32> = Vector {}
    v.push(1)
    v.push(2)
    v.push(3)
    println "pop:", v.pop()      // 3
    println "pop:", v.pop()      // 2
    println "left:", v.len()     // 1
end
\end{salam}

\begin{progout}
pop: 3
pop: 2
left: 1
\end{progout}

\subsection{The common vector operations}

The methods you will reach for most often:

\begin{center}
\begin{tabular}{@{}ll@{}}
\toprule
\textbf{Method} & \textbf{What it does} \\
\midrule
\code{v.push(x)}     & append \code{x} to the end \\
\code{v.pop()}       & remove and return the last element \\
\code{v.len()}       & the number of elements \\
\code{v.get(i)[0]}   & read the element at index \code{i} \\
\code{v.set(i, x)}   & overwrite the element at index \code{i} \\
\code{v.is\_empty()} & \code{true} if the vector has no elements \\
\code{v.cap()}       & current capacity (see the In Depth box) \\
\code{v.free()}      & release the vector's memory \\
\bottomrule
\end{tabular}
\end{center}

\begin{warn}
When you are finished with a vector, call \code{v.free()} to release the memory it
allocated. Unlike a fixed array (whose storage is reclaimed automatically when it
goes out of scope), a vector manages a growable buffer that you are responsible
for freeing. We discuss this fully in Chapter~18; for now, make \code{v.free()}
the last thing you do with a vector.
\end{warn}

\begin{deep}[In Depth: capacity versus length, and why push is cheap]
A vector keeps two numbers: its \emph{length} (how many elements you have stored)
and its \emph{capacity} (how many it could store before needing a bigger buffer).
When you push past the capacity, the vector quietly allocates a larger buffer ---
typically doubling and copies the elements across. Because the capacity grows
geometrically, the \emph{average} cost of a \code{push} is constant, even though
an occasional one triggers a copy. That is why building a vector element by
element is efficient, and why \code{v.cap()} is usually larger than
\code{v.len()}: pushing two items into a fresh vector may already reserve room
for four.
\end{deep}

\section{Arrays or vectors?}

\begin{itemize}[leftmargin=1.4em]
  \item Use a \textbf{fixed array} when you know the number of elements when you
        write the code, and it will not change coordinates, a fixed lookup
        table, the cells of a board.
  \item Use a \textbf{vector} when the count is decided at run time or changes as
        the program works reading input, accumulating results, anything
        ``however many there turn out to be.''
\end{itemize}

\section{A worked example: collecting even numbers}

This program walks the numbers 1 to 10, collects the even ones into a vector, and
reports what it found a pattern (filter into a growing collection) you will use
constantly:

\begin{salam}
func main:
    mut evens Vector<i32> = Vector {}
    for mut n = 1; n <= 10; n = n + 1:
        if n % 2 == 0:
            evens.push(n)
        end
    end
    print "found "
    print evens.len()
    println " even numbers"
    mut i i32 = 0
    while i < evens.len():
        print evens.get(i)[0]
        print " "
        i += 1
    end
    println ""
    evens.free()
end
\end{salam}

\begin{progout}
found 5 even numbers
2 4 6 8 10
\end{progout}

\section{Chapter summary}

\begin{itemize}[leftmargin=1.4em]
  \item A fixed array \code{T[N]} holds \code{N} values of type \code{T};
        elements are indexed from \code{0} to \code{N-1}.
  \item Walk an array with \code{for mut i = 0; i < N; i = i + 1}.
  \item Arrays can nest (\code{T[R][C]}) to form grids, indexed row-first.
  \item A \code{Vector<T>} grows at run time; create it with \code{Vector \{\}}.
  \item Key vector methods: \code{push}, \code{pop}, \code{len}, \code{get(i)[0]}
        to read, \code{set(i, x)} to write, \code{is\_empty}, \code{free}.
  \item Call \code{v.free()} when done with a vector.
  \item Choose arrays for fixed counts, vectors for counts that vary.
\end{itemize}

\section{Exercises}

\exercise Create an array of the five vowels (\code{char} values) and print each
on its own line with a loop.

\exercise Given \code{a i32[6] = [4, 8, 15, 16, 23, 42]}, compute and print the
average as an \code{f64}.

\exercise Build a vector of the squares $1, 4, 9, \dots, 100$ using a loop, print
its length, then print every element.

\exercise Write a program that reads a fixed array of integers and prints the
minimum and maximum, using a single loop.

\exercise Fill a \code{i32[3][3]} grid so that \code{grid[r][c] = r * 3 + c}, then
print it as three rows of three numbers.

\exercise Use a vector as a stack: push the numbers 1 to 5, then pop and print
them all, observing that they come out in reverse order. Remember to
\code{free} the vector.
